\documentclass[conference]{IEEEtran}
\usepackage{cite}
\usepackage{amsmath,amssymb,amsfonts}
\usepackage{graphicx}
\usepackage{textcomp}
\usepackage{xcolor}

\begin{document}

\title{Toward Multi-Recognizer Multi-Adversary Training for Robust Face Recognition}

\author{
\IEEEauthorblockN{Jerome Jodie Harianto}
\IEEEauthorblockA{
FYP Midterm Report\\
CUHK-Shenzhen 123010005@link.cuhk.edu.cn
}
}

\maketitle

\begin{abstract}
Manipulated face imagery poses a growing threat to biometric verification
systems because modern face synthesis and identity-targeted attacks can alter
similarity relationships in embedding space while remaining visually plausible.
This project initially focused on building a practical local framework for
running diverse face-manipulation attack families and evaluating them across
multiple recognizers. The work has since been refined toward a stronger
research question: how adversarial ensemble training can be extended more
thoroughly and more appropriately to face recognition by explicitly accounting
for adversary-family diversity, recognizer-family diversity, and
embedding-space verification behavior. Importantly, this pivot does not
discard prior work; the implemented attack backends, recognizer integrations,
and evaluation infrastructure now form the basis of the new training-oriented
study. The ArcFace training model is already integrated with the intended
ResNet-100 backbone confirmed, while the final end-to-end training pipeline
and exact training method remain under construction. The final paper will
study whether multi-recognizer multi-adversary training improves robustness to
seen and unseen manipulations while keeping clean verification performance
within an acceptable range.
\end{abstract}

\begin{IEEEkeywords}
face recognition, adversarial robustness, biometric security, face manipulation, ArcFace
\end{IEEEkeywords}

\section{Introduction}
Face recognition systems are increasingly deployed in security-sensitive
settings such as identity verification, access control, and digital forensics.
At the same time, face swapping, morphing, reenactment, and adversarial
identity attacks have become more capable and accessible. This creates a
practical security problem: a manipulated probe may remain convincing enough to
mislead a recognizer even when the image is not genuine.

An initial response is to benchmark attacks and measure how often they succeed
against existing recognizers. While this remains necessary, it is not by itself
the strongest research contribution for the current project. A more compelling
direction is recognizer-side hardening: instead of only measuring
vulnerability, the recognizer itself should be trained or fine-tuned to become
more robust to diverse manipulation families.

This midterm report summarizes the project's transition toward
\emph{multi-recognizer multi-adversary training} for robust facial embedding
learning. The central idea is that robustness may depend not only on exposure
to a strong single attacker, but also on a structured ensemble on both sides:
multiple recognizers and multiple attack mechanisms with distinct roles. This
report therefore focuses on completed groundwork, the rationale for the pivot,
and the current plan for the final paper.

\section{Background and Partial Related Work}
Modern face recognition systems commonly rely on embedding-based learning,
where facial images are mapped into a feature space designed for verification
and identification. Margin-based methods such as ArcFace are strong baselines
because they improve embedding discrimination under face-recognition objectives
\cite{arcface}. Related recognizer families such as CosFace and CurricularFace
motivate the idea that robustness should not be studied through only one
recognizer design.

Prior work has also shown that face-recognition systems are vulnerable to
diverse manipulation families, including face swapping, morphing, and
identity-targeted generative attacks \cite{simswap,advfacegan,mordiff}. These
attack families differ in how they alter the input and how well they transfer
across recognizers. As a result, a recognizer hardened against only one
manipulation family may still remain fragile under unseen or cross-family
conditions.

This motivates the refined project direction. Rather than treating attacks only
as benchmark inputs, the final paper will investigate whether diverse adversary
exposure during recognizer training produces stronger and more transferable
robustness than strongest-single-attack training alone. At this stage,
the literature review is still incomplete, but the conceptual basis is already
clear: face-recognition robustness should be studied in an embedding-centered
setting with both clean-performance and robustness tradeoffs reported
explicitly \cite{altuncu2025}.

\section{Work Completed to Date}
Substantial preparatory work has already been completed, and this progress
remains directly useful despite the pivot in paper direction. First, multiple
face-manipulation and attack-family implementations have been integrated into
the local project environment, with practical smoke-test or verification paths
established for a core subset. This means the project already has a usable
adversary pool rather than only a conceptual training plan. The implementation
notes also distinguish benchmark-faithful, approximation-grade, and
smoke-test-only paths so later experiments can be reported honestly.

Second, recognizer and evaluation infrastructure has already been prepared.
Multiple recognizers have been set up as part of the local framework, including
modern and legacy references used for cross-recognizer evaluation. This matters
because the revised paper direction depends on recognizer-specific and
cross-recognizer analysis rather than a single pooled verification result.

Third, the ArcFace training model has already been brought into the project and
the main training backbone has been fixed to ResNet-100. An earlier training
check helped confirm that recognizer training in the local environment is
practical, even though the final end-to-end training pipeline and final method
details are not yet complete. This is still meaningful progress because the
main implementation risk has shifted from backbone availability to pipeline
assembly, attack scheduling, and evaluation design.

Taken together, these completed steps show that the project is not restarting.
The current state already includes attack-family implementations, confirmed
recognizer families, evaluation scaffolding, and a usable ArcFace training
entry point. These now serve as the operational base for the final paper.

\section{Pivot in Research Direction}
The project originally leaned more heavily toward attack benchmarking and
broader evaluation infrastructure for manipulated-face security analysis. That
work was still useful because it forced the construction of a local framework
capable of running multiple attack families and multiple recognizers under
practical device constraints. However, as the project matured, it became clear
that a pure benchmark-oriented story would be less focused and less compelling
than a training-oriented contribution centered on recognizer robustness.

The revised direction is therefore a refinement rather than a reset. The
project now aims to study multi-recognizer multi-adversary training for robust
face recognition. In this framing, previously integrated attack methods become
the adversary pool used for robust training and evaluation, while the existing
recognizer setup becomes the basis for same-family and cross-family robustness
analysis. The preliminary ArcFace training test further supports the
feasibility of this pivot by showing that recognizer fine-tuning can already be
exercised in practice.

This change also sharpens the scientific question. Instead of asking only which
attacks succeed against recognizers, the project now asks which training
assembly best improves robustness: strongest single-attack training,
primary-attacker-set training, or primary-plus-surrogate attacker training
across recognizer families.

\section{Current Method Direction}
The current planned direction is to train or fine-tune a face recognizer
against a structured adversary pool and evaluate whether robustness
generalizes beyond the attacks seen during training. On the recognizer side,
the ensemble is intentionally restricted to ArcFace, CosFace, and
CurricularFace. ArcFace is already the main recipe with a confirmed ResNet-100
backbone, while CosFace and CurricularFace provide closely related but still
distinct margin-based formulations for same-family robustness comparison.

On the adversary side, the plan is organized into a primary attacker set and a
surrogate attacker set. The primary attacker set is PGD, BPFA, and DFANet,
chosen to cover white-box optimization, modern transfer pressure, and
face-recognition-specific feature-level attack pressure. The surrogate
attacker set is AdvFaceGAN, Adv-Makeup, and Greedy-DiM, chosen to cover
unrestricted dual-identity pressure, semantic makeup manipulation, and
morphing pressure. This structure is intended to keep the attack ensemble
grounded and role-based rather than treating every available attack as one flat
list.

At this stage, the final end-to-end training pipeline, exact attack sampling
policy, and final loss formulation are not finished yet. The final paper will
therefore compare conditions such as no adversarial training, strongest single
attack, primary-attacker-set training, and primary-plus-surrogate mixtures
only after the pipeline is fully assembled and validated.

\section{Next Steps}
The next project stage is to convert the completed infrastructure and confirmed
model choices into a full experimental pipeline. Since ArcFace with ResNet-100
is already fixed and the attack-family scope is already defined, the immediate
tasks are now to assemble the final training loop, finalize the attack
sampling and scheduling policy, and lock the first comparison conditions for
training: clean baseline, strongest single attacker, primary-attacker-set
training, and primary-plus-surrogate training.

After that, the project will execute reduced local smoke experiments to
validate data flow and training logic, followed by larger-scale training runs
under the final configuration where available. The literature review and
experimental sections will also be expanded substantially for the final paper.

\begin{thebibliography}{00}
\bibitem{arcface} J. Deng, J. Guo, N. Xue, and S. Zafeiriou, ``ArcFace: Additive angular margin loss for deep face recognition,'' in \emph{Proc. IEEE/CVF Conf. Comput. Vis. Pattern Recognit.}, 2019, pp. 4690--4699.

\bibitem{simswap} R. Chen, X. Chen, B. Ni, and Y. Ge, ``SimSwap: An efficient framework for high fidelity face swapping,'' in \emph{Proc. ACM Int. Conf. Multimedia}, 2020, pp. 2003--2011.

\bibitem{advfacegan} Y. Huang \emph{et al.}, ``AdvFaceGAN: A face dual-identity impersonation attack method based on generative adversarial networks,'' \emph{PeerJ Comput. Sci.}, vol. 11, 2025.

\bibitem{mordiff} S. Venkatesh \emph{et al.}, ``MorDIFF: Recognition vulnerability and attack detectability of face morphing attacks created by diffusion autoencoders,'' in \emph{Proc. IEEE Int. Joint Conf. Biometrics}, 2022.

\bibitem{altuncu2025} D. Altuncu \emph{et al.}, ``Advancing facial recognition security through adversarial attack research and standardization,'' \emph{Scientific Reports}, vol. 15, 2025.
\end{thebibliography}

\end{document}
